\documentclass[12pt,a4paper]{article}

\usepackage[a4paper,margin=1in]{geometry}
\usepackage{listings}
\usepackage{xcolor}
\usepackage{setspace}
\usepackage{titlesec}

\definecolor{codegreen}{rgb}{0,0.45,0}
\definecolor{codegray}{rgb}{0.35,0.35,0.35}
\definecolor{codepurple}{rgb}{0.45,0,0.65}
\definecolor{backcolour}{rgb}{0.97,0.97,0.97}

\lstdefinestyle{program}{
    backgroundcolor=\color{backcolour},
    commentstyle=\color{codegreen},
    keywordstyle=\color{blue},
    numberstyle=\tiny\color{codegray},
    stringstyle=\color{codepurple},
    basicstyle=\ttfamily\footnotesize,
    breaklines=true,
    keepspaces=true,
    numbers=left,
    numbersep=8pt,
    showspaces=false,
    showstringspaces=false,
    showtabs=false,
    tabsize=2,
    frame=single,
    rulecolor=\color{black!35}
}

\lstdefinelanguage{R}{
    morekeywords={function,if,else,for,while,repeat,in,next,break,
                  TRUE,FALSE,NA,NULL,library,return},
    sensitive=true
}

\lstdefinelanguage{Python}{
    morekeywords={and,as,assert,break,class,continue,def,del,elif,else,
                  except,False,finally,for,from,global,if,import,in,is,
                  lambda,None,not,or,pass,raise,return,True,try,while,
                  with,yield},
    sensitive=true
}

\lstset{style=program}
\setlength{\parindent}{0pt}
\setlength{\parskip}{6pt}
\onehalfspacing

\titleformat{\section}
  {\bfseries\fontsize{16}{19}\selectfont}
  {\thesection.}{0.5em}{}
\titleformat{\subsection}
  {\bfseries\fontsize{16}{19}\selectfont}
  {}{0em}{}

\newcommand{\labheader}[2]{%
    \noindent\textit{\fontsize{11}{13}\selectfont\textbf{EXPERIMENT NO:}} #1
    \hfill
    \textit{\fontsize{11}{13}\selectfont\textbf{DATE:}} #2
    \par\vspace{1.5cm}
}

\newcommand{\labtitle}[1]{%
    \begin{center}
        {\bfseries\fontsize{16}{19}\selectfont #1}
    \end{center}
    \vspace{0.7cm}
}

\newcommand{\labheading}[1]{%
    \subsection*{#1}
    \vspace{0.15cm}
}

\newcommand{\cycleheading}[1]{%
    \clearpage
    \begin{center}
        {\bfseries\fontsize{16}{19}\selectfont CYCLE #1}
    \end{center}
    \vspace{0.7cm}
}

\newcommand{\experiment}[8]{%
    \clearpage
    \labheader{#1}{\today}
    \labtitle{#2}

    \labheading{AIM}
    {\fontsize{12}{15}\selectfont #3\par}

    \labheading{THEORY}
    {\fontsize{12}{15}\selectfont #4\par}

    \labheading{ALGORITHM}
    {\fontsize{12}{15}\selectfont
    \begin{enumerate}
        #5
    \end{enumerate}
    }

    \labheading{PROGRAM}
    \lstinputlisting[language=#7]{#6}

    \labheading{OUTPUT}
    {\fontsize{12}{15}\selectfont #8\par}

    \labheading{RESULT}
    {\fontsize{12}{15}\selectfont The program was executed successfully and the required result was obtained.\par}
}

\begin{document}

\cycleheading{I}

\experiment{1}
{Numerical Operations}
{To perform basic numerical operations on a vector using R.}
{R provides built-in functions for statistical calculations, sorting, mathematical transformations, and vector arithmetic.}
{
    \item Create a numeric vector.
    \item Find the minimum, maximum, sum, mean, and length.
    \item Sort the vector in ascending and descending order.
    \item Calculate square roots and multiply each element by two.
}
{numerical_operations.R}
{R}
{The requested numerical statistics and transformed vectors are displayed.}

\experiment{2}
{Matrix Operations}
{To perform arithmetic and algebraic operations on matrices using R.}
{R supports element-wise arithmetic, matrix multiplication, transpose, inverse, identity matrix construction, and determinant calculation.}
{
    \item Create two matrices.
    \item Perform addition, subtraction, multiplication, and division.
    \item Perform floor division and matrix multiplication.
    \item Find the transpose, inverse, identity matrix, and determinant.
}
{matrix_operations.R}
{R}
{The results of all matrix operations are displayed.}

\experiment{3}
{Data Preprocessing}
{To handle missing values and normalize numerical attributes in the Iris dataset.}
{Data preprocessing includes identifying missing values, imputing missing values, and scaling attributes. Min--max normalization maps values to the interval \([0,1]\).}
{
    \item Load the Iris dataset.
    \item Introduce and identify missing values.
    \item Replace missing values with the mean.
    \item Normalize the numerical columns using min--max normalization.
}
{data_preprocessing_op.R}
{R}
{The cleaned and normalized Iris data frame is displayed.}

\cycleheading{II}

\experiment{4}
{Linear Regression}
{To model the relationship between weight and height using simple linear regression.}
{Simple linear regression estimates a straight-line relationship \(y=b_0+b_1x\), where \(b_0\) is the intercept and \(b_1\) is the slope.}
{
    \item Generate a random weight--height data frame.
    \item Fit a linear regression model.
    \item Display the estimated regression equation.
    \item Plot the observations and fitted regression line.
    \item Predict height for a new weight value.
}
{Linear_Regression.R}
{R}
{The regression equation, fitted plot, and prediction are produced.}

\experiment{5}
{K-Means Clustering}
{To group Iris observations into three clusters using the K-Means algorithm.}
{K-Means assigns observations to the nearest centroid and repeatedly updates centroids until the assignments converge.}
{
    \item Load the first four attributes of the Iris dataset.
    \item Set the number of clusters to three.
    \item Run K-Means with multiple random starts.
    \item Display the cluster summary and plot the observations with centroids.
}
{K_Means.R}
{R}
{The K-Means summary and cluster plot are generated.}

\experiment{6}
{KNN Classification}
{To classify breast cancer observations using the K-Nearest Neighbors algorithm.}
{KNN classifies a test observation using the majority class among its \(k\) nearest training observations. Feature normalization and a confusion matrix are used to evaluate the predictions.}
{
    \item Load and clean the breast cancer dataset.
    \item Convert feature columns to numeric values and normalize them.
    \item Split the data into training and testing sets.
    \item Apply KNN with \(k=5\).
    \item Calculate accuracy for multiple values of \(k\).
    \item Display the plots and confusion matrix.
}
{KNN.R}
{R}
{The predicted classes, accuracy plot, class distribution, and confusion matrix are displayed.}

\experiment{7}
{Market Basket Analysis}
{To find frequent itemsets and association rules from transaction data.}
{The Apriori algorithm identifies itemsets whose support is above a minimum threshold and generates rules satisfying the required confidence threshold. Lift measures the strength of an association.}
{
    \item Define the transaction database.
    \item Set the minimum support and confidence.
    \item Generate candidate itemsets and calculate support.
    \item Retain frequent itemsets.
    \item Generate rules and calculate support, confidence, and lift.
    \item Plot itemset support, rule confidence, and rule lift.
}
{apriori.R}
{R}
{Frequent itemsets, association rules, and their plots are produced.}

\cycleheading{III}

\experiment{8}
{RDD Creation Using Parallelization}
{To create an RDD from a list, inspect its partitions, and perform a map transformation.}
{An RDD is a distributed, fault-tolerant collection. The \texttt{parallelize} method distributes local data across a specified number of partitions.}
{
    \item Create a Spark context.
    \item Parallelize a list into four partitions.
    \item Display the number and contents of the partitions.
    \item Square every element using \texttt{map}.
    \item Collect and display the result.
}
{rdd_parallelize.py}
{Python}
{The number of partitions, partition contents, and squared values are displayed.}

\experiment{9}
{RDD Creation from an External File}
{To read an external text file into an RDD and count its lines.}
{The Spark \texttt{textFile} method creates an RDD whose elements are the lines of a text file. RDD actions such as \texttt{collect} and \texttt{count} retrieve results.}
{
    \item Create a Spark context.
    \item Read \texttt{input.txt} as an RDD.
    \item Collect and display all lines.
    \item Count and display the number of lines.
    \item Stop the Spark context.
}
{read_file_rdd.py}
{Python}
{The contents of the external file and the total number of lines are displayed.}

\experiment{10}
{PySpark DataFrame Operations}
{To read an Iris CSV file into a PySpark DataFrame and perform queries on it.}
{A Spark DataFrame is a distributed table with named columns. Schema inference detects data types, while DataFrame filters express relational queries.}
{
    \item Create a Spark session.
    \item Read the CSV file with a header and inferred schema.
    \item Print the schema and display rows.
    \item Apply AND, OR, and NOT filter conditions.
    \item Collect a matching row and print its fields.
}
{pyspark5.py}
{Python}
{The inferred schema, filtered rows, and selected row fields are displayed.}

\experiment{11}
{Word Count Using MapReduce}
{To count words in a text file using Spark transformations and actions.}
{The MapReduce pattern maps each word to a key--value pair \((word,1)\), then reduces values with the same key by addition.}
{
    \item Create a Spark context and read the input file.
    \item Remove punctuation and convert text to lowercase.
    \item Split lines into words.
    \item Map each word to \((word,1)\).
    \item Reduce by key and collect the word counts.
}
{word_count.py}
{Python}
{The cleaned word-count results are displayed.}

\end{document}
